\chapter{Thiết lập môi trường và kiến trúc dự án}
\label{ch:moi-truong}


\section{Công cụ cần thiết}

\subsection{Thư viện Python}

File \texttt{requirements.txt} đã được \textbf{ghim đúng phiên bản} mà báo cáo dùng (trước đây chỉ ghi tên gói, người đọc cài sau sẽ nhận bản khác). Bảng dưới ghi phiên bản đo được trong \texttt{venv} tại thời điểm chạy lại toàn bộ số liệu.

\begin{table}[H]
  \centering
  \small
  \begin{tabularx}{\textwidth}{>{\hsize=1.5\hsize}X l >{\hsize=0.5\hsize}X}
    \toprule
    \textbf{Gói} & \textbf{Phiên bản đo được} & \textbf{Mục đích} \\
    \midrule
    \texttt{llama-index-core} & 0.14.24 & Framework RAG \\
    \texttt{llama-index-embeddings-huggingface} & 0.8.0 & Nạp mô hình embedding \\
    \texttt{llama-index-llms-ollama} & 0.11.0 & Kết nối LLM qua Ollama \\
    \texttt{sentence-transformers} & 6.0.1 & Chạy embedding/reranker \\
    \texttt{torch} (cài kèm) & 2.14.0 (CUDA 13.0) & Chạy mô hình trên GPU \\
    \texttt{sentencepiece}, \texttt{protobuf} & 0.2.2, không đo & Phụ thuộc của tokenizer \\
    \texttt{pytesseract} & 0.3.13 & Gọi Tesseract từ Python \\
    \texttt{pdf2image} & 1.17.0 & Render PDF thành ảnh \\
    \texttt{httpx} & 0.28.1 & Gọi HTTP tới Ollama và model server \\
    \texttt{ollama} & 0.6.2 & Client Python của Ollama (bắt lỗi trong \texttt{app/app.py}) \\
    \texttt{gradio} & 6.27.0 & Giao diện web chạy thử (\texttt{app/app.py}) \\
    \texttt{python-dotenv} & 1.2.3 & Đọc biến môi trường từ \texttt{.env} \\
    \bottomrule
  \end{tabularx}
  \caption{Thư viện Python trong \texttt{requirements.txt}.}
\end{table}

\subsection{Phụ thuộc hệ thống}

\begin{itemize}
  \item \textbf{Tesseract OCR} 5.5.3, đã xác nhận có gói ngôn ngữ \texttt{vie}.
  \item \textbf{poppler-utils} (\texttt{pdftoppm} 26.01.0), cần cho \texttt{pdf2image}.
  \item \textbf{Ollama} để chạy LLM (script \texttt{pull\_models.sh} kiểm tra và báo lỗi nếu chưa cài).
  \item Python 3.14.6 trong \texttt{venv}.
\end{itemize}

\subsection{LLM chạy qua Ollama}

Để dữ liệu không đi qua dịch vụ AI bên thứ ba, LLM được chạy qua \textbf{Ollama}, một trình phục vụ mô hình mở với API HTTP (mặc định tại \texttt{http://localhost:11434}); ở profile máy chủ, Ollama chạy trên một máy GPU thuê riêng và địa chỉ được đọc từ biến \texttt{OLLAMA\_BASE\_URL}. Đồ án dùng họ mô hình Qwen3: \texttt{qwen3:4b} cho profile chạy cục bộ và \texttt{qwen3:8b} cho profile máy chủ. LlamaIndex kết nối tới Ollama~\cite{ollama} qua gói \texttt{llama-index-llms-ollama}; Qwen3 là dòng mô hình mở của Alibaba, bản 4B và 8B đủ nhỏ để chạy trên máy cá nhân~\cite{qwen3}. Thời gian chờ mặc định trong cấu hình là 180 giây do mô hình chạy trên phần cứng cá nhân có thể trả lời chậm. Qwen3 có chế độ ``suy nghĩ'' (sinh chuỗi lập luận trước câu trả lời); đồ án tắt chế độ này (\texttt{thinking=False}) để câu trả lời ngắn và nhanh.


\subsection{Phần cứng và bộ nhớ GPU}
\label{sec:phan-cung}

Máy phát triển dùng card rời NVIDIA RTX A2000 Laptop (4 GB VRAM, khoảng 3{,}7 GB khả dụng cho PyTorch). Kiểm tra bằng cách nạp mô hình qua đúng hàm của dự án và in thiết bị của các tham số cho thấy cả mô hình embedding lẫn reranker đều nằm trên \texttt{cuda:0}. LLM không chạy trên máy này: ở profile \texttt{server}, Ollama chạy trên máy GPU thuê từ dịch vụ ckey.vn, được quản lý bằng \texttt{scripts/manager.py}.

Ở bản báo cáo trước, profile \texttt{server} \textbf{không nạp đồng thời} được embedding và reranker trên card 4 GB: tiến trình đã chiếm khoảng 3{,}5 GB và báo \texttt{CUDA out of memory} khi nạp mô hình thứ hai, nên phép đo phải đẩy embedding lên CPU. Bản hiện tại đã gộp thay đổi fp16: embedding còn khoảng 1{,}1 GB, cả hai mô hình cùng nằm trên GPU tổng khoảng 2{,}2 GB (xem \texttt{src/index\_builder.py} và \texttt{make\_local\_reranker()} trong \texttt{src/retriever.py}). Thêm nữa, \texttt{resolve\_device()} tự hạ cấp về CPU khi máy không có CUDA, nên cùng mã nguồn chạy được cả trên máy GPU, máy CPU và CI. Các số liệu ở Chương~\ref{ch:danh-gia} được đo lại trên \textbf{CPU} vì môi trường chạy lại không có driver NVIDIA; đây là lý do thời gian truy hồi trong báo cáo lớn hơn con số đo trên GPU ở bản trước.

\section{Cài đặt và Hello World}

\subsection{Các bước cài đặt}

Trình tự đã dùng trong dự án:

\begin{enumerate}
  \item Tạo và kích hoạt môi trường ảo Python (\texttt{venv}).
  \item Cài các gói Python: \texttt{pip install -r requirements.txt}.
  \item Cài Tesseract (kèm gói ngôn ngữ \texttt{vie}) và poppler-utils bằng trình quản lý gói của hệ điều hành.
  \item Cài Ollama, chạy \texttt{ollama serve}, rồi tải mô hình bằng \texttt{scripts/pull\_models.sh [local|server]}.
\end{enumerate}

\subsection{Hello World}

Script \texttt{scripts/hello\_world.py} kiểm tra nhanh rằng LlamaIndex cài đúng: tạo một \texttt{Document}, dựng \texttt{SummaryIndex} và đếm số Document trong \texttt{docstore}. Script này không cần LLM hay mô hình embedding nên chạy được ngay sau khi cài thư viện.

\begin{lstlisting}[language=Python, caption={Tạo \texttt{Document} và \texttt{SummaryIndex} với LlamaIndex (\texttt{scripts/hello\_world.py}).}, label={lst:hello-world}]
from llama_index.core import Document, SummaryIndex

doc = Document(text = "(*@Điều 15. Sinh viên nghỉ học quá 20\% số tiết sẽ không đủ điều kiện dự thi.@*)")
print("(*@Document tạo thành công!@*)")
print("(*@Nội dung:@*)", doc.text)
print("(*@ID tự sinh:@*)", doc.doc_id)

index = SummaryIndex.from_documents([doc])
print("\n(*@Index build thành công!@*)")
print("(*@Số Document trong Index:@*)", len(index.docstore.docs))
\end{lstlisting}

Kết quả chạy thật:

\begin{lstlisting}[caption={Kết quả chạy \texttt{hello\_world.py}.}, label={lst:hello-world-output}]
(*@Document tạo thành công!@*)
(*@Nội dung: Điều 15. Sinh viên nghỉ học quá 20\% số tiết sẽ không đủ điều kiện dự thi.@*)
(*@ID tự sinh: 3dc12976-3b4e-4711-8798-c7d0237d98a9@*)

(*@Index build thành công!@*)
(*@Số Document trong Index: 1@*)
\end{lstlisting}

Ý nghĩa: \texttt{doc\_id} được LlamaIndex tự sinh dạng UUID, và \texttt{SummaryIndex} lưu Document vào \texttt{docstore}. Câu ví dụ ``Điều 15'' trong script chỉ là văn bản minh họa tự viết, không phải nội dung trích từ quy chế.

\section{Minh chứng kiến trúc}

\subsection{Cấu trúc thư mục dự án}

\begin{lstlisting}[caption={Cấu trúc thư mục dự án.}, label={lst:cau-truc-thu-muc}]
.
|-- app/            # Giao dien Gradio chay thu (app/app.py)
|-- data/
|   |-- raw/        # PDF goc (khong commit)
|   |-- ocr/        # Tang 1: Markdown tho sau OCR, noi sua tay
|   `-- processed/  # Tang 2: Markdown co heading, sinh lai tu tang 1
|-- eval/
|   |-- evaluate.py              # Recall@k, MRR, nDCG, CI bootstrap, McNemar
|   |-- baselines.py             # baseline BM25 / khong xep hang
|   |-- evaluate_answers.py      # cham cau tra loi bang evaluator cua LlamaIndex
|   |-- metrics.py               # do do + kiem dinh thong ke
|   |-- cau_hoi_danh_gia.json    # 13 cau co dap an chuan + nhan pham vi
|   |-- test_questions_01.json   # 31 cau hoi
|   |-- test_questions_02.json   # 10 cau hoi ve QD sua doi 3183/2025
|   `-- results/                 # ket qua JSON + Markdown
|-- notebooks/      # Chay thu pipeline, chay tren GPU thue
|-- scripts/        # hello_world.py, test_ocr.py, check_model_memory.py,
|                   # pull_models.sh, manager.py
|-- src/
|   |-- corpus.py          # Danh sach van ban nap vao index
|   |-- config.py          # RAGConfig + 2 profile + resolve_device()
|   |-- processed_loader.py# Markdown tang 2 -> TextNode
|   |-- bm25.py            # BM25 thuan Python cho hybrid search
|   |-- index_builder.py   # Node -> VectorStoreIndex + fingerprint
|   |-- retriever.py       # vector / +BM25 / +rerank / nguong tuong dong
|   |-- query_engine.py    # Prompt + LLM (temperature=0, seed) + nguon
|   |-- remote_models.py   # Client goi model server qua HTTP
|   |-- model_server.py    # FastAPI giu embedding + reranker tren GPU
|   |-- ingestion.py       # BAN CU: OCR theo trang (khong con dung)
|   |-- node_parser.py     # BAN CU: cat theo regex (khong con dung)
|   `-- pipeline.py        # Hoi dap qua dong lenh
|-- tests/              # 26 kiem thu offline (unittest)
|-- storage/            # storage/<profile>/<corpus>/ (khong commit)
|-- Makefile            # make test / eval / baseline / eval-answers / report
`-- requirements.txt    # da ghim phien ban
\end{lstlisting}

Mọi script chạy từ thư mục gốc của repository, ví dụ \texttt{RAG\_PROFILE=server python src/pipeline.py}. Các module trong \texttt{src/} vừa là thư viện (được module sau \texttt{import}) vừa có khối \texttt{\_\_main\_\_} để chạy thử riêng từng giai đoạn.

\subsection{Cấu hình tập trung và hai profile}

Toàn bộ tham số RAG nằm trong \texttt{src/config.py} dưới dạng một \texttt{dataclass} \texttt{RAGConfig}. Có hai profile sẵn: \texttt{local} (máy cá nhân, mô hình nhỏ) và \texttt{server} (mô hình lớn hơn, có reranker, LLM chạy trên máy GPU thuê). Profile được chọn bằng biến môi trường \texttt{RAG\_PROFILE} (đọc cả từ file \texttt{.env}), mặc định là \texttt{local}; giá trị không hợp lệ làm hàm \texttt{get\_config()} ném \texttt{ValueError}.

\begin{lstlisting}[language=Python, caption={Khai báo \texttt{dataclass RAGConfig} trong \texttt{src/config.py}.}, label={lst:ragconfig}]
@dataclass
class RAGConfig:
    embed_model_name: str
    embed_device: str  # "cpu" or "cuda"
    llm_model_name: str
    ollama_base_url: str
    request_timeout: float = 180.0
    temperature: float = 0.0        # 0 = tat dinh (sua loi tra loi khong on dinh)
    seed: int = 42
    reranker_model_name: str | None = None
    reranker_top_n: int = 3
    use_rerank: bool = False
    similarity_top_k: int = 10
    similarity_cutoff: float | None = None   # nguong tuong dong (None = tat)
    use_hybrid: bool = False                 # BM25 + vector qua QueryFusionRetriever
    bm25_top_k: int = 10
    fusion_mode: str = "reciprocal_rerank"
    response_mode: str = "compact"
    streaming: bool = True
    persist_dir: str = ""           # trong = storage/<profile>/<corpus>
\end{lstlisting}

\begin{table}[H]
  \centering
  \small
  \begin{tabularx}{\textwidth}{l X X}
    \toprule
    \textbf{Tham số} & \textbf{Profile \texttt{local}} & \textbf{Profile \texttt{server}} \\
    \midrule
    Embedding & \texttt{bkai-foundation-models/ vietnamese-bi-encoder} (768 chiều) & \texttt{AITeamVN/Vietnamese\_Embedding} (1024 chiều) \\
    Thiết bị embedding & \texttt{cuda} & \texttt{cuda} \\
    LLM (Ollama) & \texttt{qwen3:4b}, \texttt{localhost:11434} & \texttt{qwen3:8b}, địa chỉ lấy từ \texttt{OLLAMA\_BASE\_URL} \\
    Reranker & không dùng & \texttt{AITeamVN/Vietnamese\_Reranker}, giữ top 3 \\
    Retriever & top-10 & top-10 \\
    \bottomrule
  \end{tabularx}
  \caption{Hai profile cấu hình trong \texttt{src/config.py}.}
\end{table}

Chỉ mục được lưu tách theo profile và tự build lại khi cấu hình đổi (Chương~\ref{ch:persist}).

Nguồn dữ liệu thực tế được khai báo ở \emph{một} chỗ: \texttt{src/corpus.py} (\texttt{CORPUS\_DIR} và dict \texttt{SGU\_HIEU\_LUC}). Trường \texttt{corpus\_sources} cũ cùng hai file PDF không tồn tại đã được gỡ bỏ, nên không còn tình trạng cấu hình trỏ vào file không có thật.

\subsection{Các script hỗ trợ}

\begin{itemize}
  \item \texttt{scripts/test\_ocr.py}: render toàn bộ một PDF và OCR ra file Markdown trong \texttt{data/ocr/}, kèm số trang và tổng số ký tự. Đây là script tạo ra bản OCR \texttt{data/ocr/4.\ QuyCheDaoTaoDHSG 2021.md} mà \texttt{index\_builder.py} dùng làm đầu vào.
  \item \texttt{scripts/check\_model\_memory.py}: đo RAM/VRAM khi nạp một mô hình embedding hoặc reranker (nhận \texttt{repo\_id}, \texttt{--type}, \texttt{--device}), giúp chọn mô hình vừa với phần cứng.
  \item \texttt{scripts/pull\_models.sh}: tải LLM qua Ollama và tải embedding/reranker từ HuggingFace theo profile, đồng bộ với \texttt{config.py}; dừng với thông báo lỗi nếu thiếu \texttt{ollama}, daemon chưa chạy hoặc thiếu \texttt{huggingface\_hub}.
  \item \texttt{scripts/manager.py}: công cụ dòng lệnh gom các thao tác với dịch vụ GPU thuê (xem, thuê, khởi động lại, xóa máy; xem số dư), gọi thử Ollama và chạy một server embedding.
  \item \texttt{scripts/setup\_server\_v2.sh}, \texttt{scripts/colab\_server.py}: dựng phần chạy trên GPU
        (Ollama + model server embedding/reranker) trên Google Colab hoặc máy thuê, mở URL công khai
        qua Cloudflare quick tunnel và in sẵn các dòng \texttt{OLLAMA\_BASE\_URL}, \texttt{MODEL\_SERVER\_URL}
        để dán vào \texttt{.env} ở máy local. Thư viện của model server được ghim cùng phiên bản với
        \texttt{requirements.txt}, để vector sinh trên Colab khớp với chỉ mục đã build ở local. Nhờ vậy
        phần đánh giá câu trả lời --- vốn cần LLM --- chạy được mà không cần máy có GPU.
  \item \texttt{eval/run\_batch\_cau\_hoi.py}: chạy hàng loạt danh sách câu hỏi trong
        \texttt{eval/results/cau-hoi01.md} qua query engine, ghi câu trả lời kèm nguồn ra
        \texttt{cau-hoi01-result.md} để chấm tay (Mục~\ref{sec:chay-lai-56-cau}).
  \item \texttt{Makefile}: gom \texttt{make test} (26 kiểm thử offline), \texttt{make eval} (đo 4 cấu hình truy hồi), \texttt{make baseline}, \texttt{make eval-answers} và \texttt{make report}.
  \item \texttt{.github/workflows/ci.yml}: CI cài đúng phần tối thiểu (không torch, không model) và chạy toàn bộ kiểm thử offline trên mỗi lần push.
\end{itemize}


Chuỗi xử lý của đồ án gồm sáu giai đoạn. Hình \ref{fig:pipeline} thể hiện chuỗi này và file mã nguồn tương ứng; cả sáu giai đoạn đều đã có mã chạy được. Pha chuẩn bị (1--3) chạy một lần và lưu chỉ mục xuống đĩa; pha truy vấn (4--5) chạy mỗi câu hỏi; giai đoạn 6 dùng lại retriever của giai đoạn 4.

\begin{figure}[H]
  \centering
  \begin{tikzpicture}[
      node distance=6mm,
      done/.style={rectangle, draw, rounded corners, minimum width=2.6cm, minimum height=1.1cm, align=center, font=\small},
      arr/.style={-{Stealth}, thick}
    ]
    \node[done] (n1) {1. Dữ liệu\\ \scriptsize extract\_to\_md.py $\to$ clean\_md.py};
    \node[done, right=of n1] (n2) {2. Node\\ \scriptsize processed\_loader.py};
    \node[done, right=of n2] (n3) {3. Index\\ \scriptsize index\_builder.py};
    \node[done, below=8mm of n1] (n4) {4. Retriever\\ \scriptsize retriever.py};
    \node[done, right=of n4] (n5) {5. Synthesis\\ \scriptsize query\_engine.py};
    \node[done, right=of n5] (n6) {6. Evaluation\\ \scriptsize eval/evaluate.py};
    \draw[arr] (n1) -- (n2);
    \draw[arr] (n2) -- (n3);
    \draw[arr] (n3.south) -- ++(0,-3mm) -| ([yshift=0mm]n4.north);
    \draw[arr] (n4) -- (n5);
    \draw[arr, dashed] (n4.south) -- ++(0,-4mm) -| (n6.south);
  \end{tikzpicture}
  \caption{Sáu giai đoạn của pipeline RAG và file mã nguồn tương ứng; mũi tên nét đứt: bộ đánh giá gọi trực tiếp retriever.}
  \label{fig:pipeline}
\end{figure}

\begin{table}[H]
  \centering
  \small
  \begin{tabularx}{\textwidth}{l l X l}
    \toprule
    \textbf{Giai đoạn} & \textbf{Mức} & \textbf{File} & \textbf{Chạy được} \\
    \midrule
    1. Dữ liệu & Bắt buộc lõi & \texttt{scripts/extract\_to\_md.py}, \texttt{scripts/clean\_md.py} & Có \\
    2. Node & Bắt buộc lõi & \texttt{src/processed\_loader.py} & Có \\
    3. Index & Bắt buộc lõi & \texttt{src/index\_builder.py} & Có \\
    4. Retriever & Bắt buộc lõi & \texttt{src/retriever.py} & Có \\
    5. Response Synthesis & Bắt buộc lõi & \texttt{src/query\_engine.py}, \texttt{src/pipeline.py} & Có \\
    6. Evaluation & Bắt buộc lõi & \texttt{eval/evaluate.py}, \texttt{eval/baselines.py}, \texttt{eval/evaluate\_answers.py} & Có (truy hồi: tự động; câu trả lời: cần Ollama) \\
    Hybrid BM25 & Chọn thêm & \texttt{src/bm25.py}, \texttt{src/retriever.py} & Có (mặc định tắt) \\
    Rerank & Chọn thêm & \texttt{src/retriever.py} & Có (profile \texttt{server}, mặc định tắt) \\
    Model server & Chọn thêm & \texttt{src/model\_server.py}, \texttt{src/remote\_models.py} & Có \\
    Kiểm thử & Bắt buộc lõi & \texttt{tests/}, \texttt{Makefile}, CI & Có (26 kiểm thử) \\
    Giao diện web & Chọn thêm & \texttt{app/app.py} & Có (chạy thử) \\
    \bottomrule
  \end{tabularx}
  \caption{Trạng thái mã nguồn của từng giai đoạn.}
\end{table}

\section{Kiểm chứng môi trường}
\label{sec:kiem-chung-moi-truong}

\paragraph{Chỉ mục khớp với loader.} Chạy \texttt{src/index\_builder.py} (profile \texttt{server}, CPU,
ngày 09/10):

\begin{lstlisting}[caption={Kiểm chứng số Node trong chỉ mục.}, label={lst:kiem-chung-index}]
[INFO] Load Index từ storage/server/01_quy_che_dao_tao/1_quy_che_dao_tao_dai_hoc
Index: storage/server/01_quy_che_dao_tao/1_quy_che_dao_tao_dai_hoc | số node trong Index: 54 | từ processed_loader: 54 | khớp: True
\end{lstlisting}

\paragraph{Kết nối Ollama.} Một câu trả lời thật trong nhật ký Gradio (\texttt{eval/results/questions\_log.jsonl},
ngày 08/10, \texttt{qwen3:8b}) cho câu ``Sinh viên được rút học phần trong thời gian nào?'':

\begin{lstlisting}[caption={Câu trả lời mẫu qua Ollama.}, label={lst:ollama-mau}]
Theo quy định tại Điều 7 — Quy chế đào tạo trình độ đại học SGU 2021 (sửa đổi bởi
QĐ 3183/2025), việc rút bớt học phần trong khối lượng học tập đã đăng ký được thực
hiện trong 03 tuần đầu tiên của học kì.
\end{lstlisting}

Câu trả lời khớp khoản 4 mục sửa đổi Điều 7 của QĐ 3183/2025.
